\chapter{Pulizia statistica e imputazione}
\label{sec:statistical-cleaning-and-imputation}

Il trattamento statistico descritto in questo capitolo si applica ai dati che superano la selezione per posizione ODV della \cref{sec:sales-position-selection}. La coerenza dei bilanci di massa e dei tempi non esclude la presenza di quantità o durate insolitamente alte o basse rispetto ad altre osservazioni confrontabili. Possono inoltre rimanere durate pari a zero nelle macrofasi per cui occorre rappresentare un tempo di lavorazione positivo. Si individuano questi casi, si stimano valori sostitutivi e se ne verifica la compatibilità con il ciclo produttivo.

\section{Obiettivi e criteri generali}
\label{sec:statistical-treatment-principles}

% Current implementations in StrategicPlannerBeanTech:
% data/statistical_mass.py, mass/outlier_uncertainty.py;
% data/statistical_duration.py, duration/statistical_audit.py;
% data/zero_duration.py and data/expected_duration.py.
% General principles only here. Definitions and estimation details follow
% locally, before the corresponding equations, tables and figures.
Un'incoerenza logica viola una delle condizioni adottate per ricostruire il ciclo. Un outlier statistico è un valore estremo rispetto alla distribuzione di riferimento, ma non necessariamente errato. La segnalazione individua quindi un caso da esaminare prima di un'eventuale sostituzione. Il confronto avviene con osservazioni della stessa grandezza, raggruppate secondo i criteri specificati nelle sezioni seguenti.

Anche i valori mancanti e gli zeri richiedono interpretazioni distinte. Un valore mancante indica che la misura non è disponibile. Uno zero è invece una misura presente, da interpretare in base alla variabile e alla macrofase. In particolare, le durate pari a zero vengono analizzate separatamente dalle durate positive anomale e imputate soltanto per le macrofasi per cui è assunta una durata di lavorazione positiva.

Le procedure correggono i valori senza eliminare ulteriori ODP o occorrenze.

\clearpage
\section{Trattamento statistico delle masse}
\label{sec:statistical-mass-treatment}

Il trattamento delle masse analizza le quantità iniziali, gli scarti e le quantità finali degli ODP. Dopo aver individuato i casi da correggere, le quantità vengono aggiornate insieme per mantenere i bilanci coerenti. Restano invariati il materiale pronto registrato e i collegamenti del percorso selezionato.

\subsection{Grandezze analizzate e campione di riferimento}
\label{sec:mass-statistical-variables}

% Sources: mass/statistical_audit.py, mass_views; current pre-imputation frames.
% Pre-imputation values restored in memory from current mass/duration audits.
% Point thresholds verified against the stored audit to relative tolerance 1e-10.
% Evidence and source hashes: tables/data/statistical-mass-summary.json.
% Figures: scripts/plot_statistical_mass.py and figures/data/statistical-mass-*.csv.
Si analizzano due rapporti: la resa di ciascun ODP e lo scarto relativo alla quantità lavorata in ogni occorrenza. La resa è il rapporto tra massa finale e massa iniziale. Lo scarto relativo è il rapporto tra lo scarto e la quantità in ingresso all'occorrenza.

L'indice $r(k)$ identifica l'occorrenza di \texttt{TRENOLUNA} dell'ODP $k$, mentre $q_{r(k)}$ è la sua quantità in ingresso. Indichiamo con $m_k^{\mathrm{fin}}$ la massa del nodo finale dell'ODP $k$. Per l'occorrenza $i$, $q_i$ è la quantità in ingresso e $s_i$ lo scarto. Tutte queste quantità sono espresse in chilogrammi. La resa $y_k$ e lo scarto relativo $\rho_i$ sono quindi
\begin{equation}
\label{eq:statistical-mass-ratios}
y_k=\frac{m_k^{\mathrm{fin}}}{q_{r(k)}},
\qquad
\rho_i=\frac{s_i}{q_i}.
\end{equation}
Negli ODP con completamenti impliciti si usa la massa del nodo finale, che raccoglie il materiale dei diversi rami. Le rese vengono confrontate tra tutti gli ODP. Gli scarti relativi vengono confrontati tra occorrenze della stessa macrofase.

L'analisi riguarda \num{4860} ODP, associati a \num{4455} posizioni ODV e \num{2103} ODV distinti. Per lo scarto si considerano \num{11591} occorrenze materiali, escludendo \texttt{HPT} e gli stati di materiale pronto. Di queste, \num{4812} hanno scarto nullo ($\rho_i=0$), 337 hanno scarto totale ($\rho_i=1$) e \num{6442} hanno scarto parziale ($0<\rho_i<1$). Il confronto statistico degli scarti riguarda i casi parziali. Gli scarti nulli e totali sono trattati come casi distinti.

\subsection{Individuazione dei candidati e incertezza delle soglie}
\label{sec:mass-outlier-detection}

% Implemented criterion: MAD with outer bootstrap confidence limits.
% Source: mass/outlier_uncertainty.py and statistical_mass_thresholds audit.
Per individuare i valori estremi si calcolano una soglia inferiore e una superiore, sia per la resa sia per lo scarto relativo alla quantità lavorata di ciascuna macrofase. Si usa la trasformazione logit per i rapporti $p_j$ compresi tra zero e uno, con $j$ che identifica un'osservazione. La funzione logistica $g$ riporta i valori alla scala originale:
\begin{equation}
\label{eq:mass-logit-transform}
z_j=\log\left(\frac{p_j}{1-p_j}\right),
\qquad
p_j=g(z_j)=\frac{1}{1+\exp(-z_j)},
\qquad 0<p_j<1.
\end{equation}

Il criterio MAD usa la mediana $m$ dei valori $z_j$ e la deviazione assoluta mediana, $\operatorname{MAD}=\operatorname{mediana}_j(\lvert z_j-m\rvert)$. Le soglie delimitano l'intervallo $B(m,3.5\cdot1.4826\cdot\operatorname{MAD})$, dove $B(m,r)=[m-r,m+r]$.
Il fattore \num{1.4826} rende la MAD confrontabile con la deviazione standard usando la distribuzione normale come riferimento \parencite{nist2016mad}.

Il calcolo delle soglie richiede almeno 20 rapporti e una MAD positiva, sia nei dati sia nelle singole repliche bootstrap. Con riferimenti provenienti da almeno due ODV, si eseguono \num{2000} ricampionamenti, mantenendo insieme le osservazioni dello stesso ODV. Per ogni soglia si calcola un intervallo di confidenza al 95\% (IC95) dalle repliche valide. Il criterio MAD + IC95 segnala un valore se è sotto la soglia inferiore e l'estremo inferiore del suo IC95, oppure sopra la soglia superiore e l'estremo superiore del suo IC95.

La \cref{tab:mass-outlier-thresholds} riporta i limiti in percentuale. La resa pari al 100\% di un ODP è esclusa dal logit, lasciando \num{4859} osservazioni. Per lo scarto parziale sono valutabili \num{6431} occorrenze su \num{6442}: le altre 11 appartengono a macrofasi con riferimenti insufficienti.

\begin{table}[htbp]
\centering
\small
\sisetup{group-digits=integer}
\begin{tabularx}{\linewidth}{@{}Xrrrrr@{}}
\toprule
\textbf{Ambito} & $n$ & \textbf{Limite inf. [\%]} & \textbf{Limite sup. [\%]} & \textbf{Segnalati} & \textbf{Repliche} \\
\midrule
Resa globale & \num{4859} & \num{84.540} & \num{98.484} & \num{76} & \num{2000} \\
\texttt{TRENOLUNA} & \num{4860} & \num{1.509} & \num{12.683} & \num{77} & \num{2000} \\
\texttt{BANCHI} & \num{979} & \num{0.449} & \num{86.109} & \num{4} & \num{2000} \\
\texttt{CIRCO} & \num{38} & \num{0.01623} & \num{71.089} & \num{0} & \num{1999} \\
\texttt{CONF} & \num{34} & \num{0.00473} & \num{99.429} & \num{0} & \num{1985} \\
\texttt{RADDR} & \num{68} & \num{0.01632} & \num{28.155} & \num{0} & \num{2000} \\
\texttt{TAG} & \num{419} & \num{0.00902} & \num{95.707} & \num{0} & \num{2000} \\
\texttt{TRATT\_LUNA} & \num{33} & \num{0.00994} & \num{43.106} & \num{0} & \num{1942} \\
\bottomrule
\end{tabularx}
\caption{Limiti MAD + IC95 sulla scala percentuale. $n$ è il numero di osservazioni usate: ODP per la resa, occorrenze con scarto relativo strettamente tra zero e uno per le macrofasi. «Segnalati» conta le osservazioni oltre i limiti. «Repliche» indica quante delle \num{2000} repliche bootstrap hanno almeno 20 rapporti e MAD positiva. Sono mostrati gli ambiti valutabili.}
\label{tab:mass-outlier-thresholds}
\end{table}


Sono segnalati 76 ODP per resa bassa e 81 per almeno uno scarto relativo anomalo. Poiché 56 ODP compaiono in entrambi gli insiemi, i candidati sono 101, pari al \num{2.08}\% del campione. Le segnalazioni locali riguardano 77 occorrenze \texttt{TRENOLUNA} e quattro \texttt{BANCHI}. La \cref{fig:mass-outlier-distributions} mostra le distribuzioni e le soglie prima dell'imputazione.

\begin{figure}[H]
\centering
\reportgraphic[height=0.42\textheight]{statistical-mass-distributions.pdf}
\caption{Istogrammi prima dell'imputazione. Le classi hanno uguale ampiezza sulla scala logit, mentre le etichette orizzontali riportano percentuali. Le barre gialle identificano le osservazioni segnalate da MAD + IC95. Le linee tratteggiate mostrano le soglie MAD e le bande i rispettivi IC95. Negli istogrammi dello scarto relativo sono esclusi i casi di scarto nullo e totale.}
\label{fig:mass-outlier-distributions}
\end{figure}

\clearpage
\subsection{Stima dei rapporti da usare per l'imputazione}
\label{sec:mass-ratio-estimation}

% Implementation: mass/imputation_models.py, fit_imputers and RatioImputer.
% Current cached bundle: outputs/models/statistical_mass/
% 5b55c4f9adb9bb8324d89ef165684cab80815cde56ffd036380bf1ccd423fada.joblib.
% Held-out verification computed without retraining; source hashes in summary.
Una resa o uno scarto relativo anomalo non indicano da soli quale quantità in chilogrammi correggere. Si stimano prima rapporti plausibili, escludendo dai dati di riferimento tutti i 90 ODV contenenti un ODP candidato. I \num{2013} ODV rimanenti sono divisi in \num{1207} per addestrare i modelli, 403 per calibrare gli intervalli e 403 per verificarne il comportamento. Posizioni e ODP dello stesso ODV rimangono nello stesso gruppo.

Si usano due regressioni Ridge sulla scala logit, una per la resa e l'altra per lo scarto relativo parziale. Le variabili esplicative sono quantità richiesta della posizione ODV, prima dimensione, lunghezza richiesta, numero di ODP della posizione, prodotto, forma, finitura, stato fisico e marca dell'acciaio. Il modello dello scarto relativo include anche la macrofase.

Per quantità richiesta, dimensione e lunghezza, ciascun valore $x$ viene trasformato con \mbox{$\log(1+x)$}. I valori numerici mancanti sono sostituiti con la mediana dell'addestramento. A ciascuna variabile numerica si sottrae la media e si divide il risultato per la deviazione standard, entrambe calcolate sull'addestramento. Le variabili categoriali sono rappresentate da indicatori zero oppure uno, con una categoria \texttt{MISSING} per i valori mancanti. Una categoria mai vista nell'addestramento azzera gli indicatori della sola variabile interessata, lasciando gli altri contributi alla stima.

Indichiamo con $\bm{x}_j$ il vettore risultante per l'osservazione $j$, con $\bm{\beta}$ i coefficienti e con $\beta_0$ l'intercetta. L'addestramento minimizza
\begin{equation}
\label{eq:mass-ridge-objective}
\min_{\beta_0,\bm{\beta}}\;
\sum_{j\in\mathcal{T}}\left(z_j-\beta_0-\bm{x}_j^{\mathsf{T}}\bm{\beta}\right)^2
+\lambda\lVert\bm{\beta}\rVert_2^2,
\qquad \lambda=10,
\end{equation}
dove $\mathcal{T}$ è l'insieme delle osservazioni di addestramento. La penalità limita l'ampiezza dei coefficienti, senza penalizzare l'intercetta. Le stime sono
\begin{equation}
\label{eq:mass-ridge-estimation}
\widehat{z}_j=\beta_0+\bm{x}_j^{\mathsf{T}}\bm{\beta},
\qquad
\widehat{p}_j=g(\widehat{z}_j),
\end{equation}
con $g$ definita nella \cref{eq:mass-logit-transform}.

Nel gruppo di calibrazione si calcolano i residui $\varepsilon_j=z_j-\widehat{z}_j$. Per lo scarto relativo si usano quelli della stessa macrofase se sono almeno 30, altrimenti quelli complessivi del modello. Detti $c_{0.05}$ e $c_{0.95}$ i percentili 5 e 95 dei residui così selezionati, l'intervallo empirico è
\begin{equation}
\label{eq:mass-ratio-interval}
\left[g\left(\widehat{z}_j+c_{0.05}\right),\,
      g\left(\widehat{z}_j+c_{0.95}\right)\right].
\end{equation}
L'intervallo ha livello nominale del 90\%, poiché usa i percentili 5 e 95 dei residui di calibrazione. L'IC95 della sezione precedente riguarda invece l'incertezza delle soglie per gli outlier. La suddivisione circa 60/20/20, $\lambda=10$ e il minimo di 30 residui sono impostazioni fisse della procedura, senza una ricerca dei valori migliori.

La \cref{tab:mass-imputer-validation} riporta i risultati nel gruppo di verifica. Il MAE è l'errore assoluto medio tra rapporti osservati e stimati, espresso in punti percentuali (p.p.). La copertura è la percentuale di osservazioni comprese nel proprio intervallo. Per lo scarto relativo è dell'\num{87.87}\%, inferiore di \num{2.13} punti percentuali al livello nominale. La copertura complessiva non garantisce lo stesso risultato per ogni macrofase. Queste misure valutano le previsioni sui dati di riferimento, non l'accuratezza delle correzioni dei 101 ODP candidati, per i quali mancano valori reali verificati.

\begin{table}[H]
\centering
\small
\begin{tabularx}{\linewidth}{@{}Xrrr@{}}
\toprule
\textbf{Variabile} & \textbf{Osservazioni} & \textbf{MAE [p.p.]} & \textbf{Copertura [\%]} \\
\midrule
Resa globale & \num{911} & \num{1.29} & \num{90.12} \\
Scarto relativo parziale & \num{1195} & \num{2.35} & \num{87.87} \\
\bottomrule
\end{tabularx}
\caption{Risultati nel gruppo di verifica. Le osservazioni sono rapporti di resa oppure scarti relativi parziali. Il MAE è espresso in punti percentuali. La copertura è la percentuale di osservazioni comprese negli intervalli empirici nominali al 90\%.}
\label{tab:mass-imputer-validation}
\end{table}


Nella riconciliazione si usano le stime della resa di ogni ODP candidato e dello scarto relativo di tutte le sue occorrenze con scarto parziale, anche se non segnalate. Per le occorrenze con scarto nullo o totale non si usa una stima dello scarto relativo.

\subsection{Riconciliazione delle quantità in chilogrammi}
\label{sec:mass-constrained-imputation}

% Implementations: mass/constrained_imputation.py, reconcile_family;
% mass/plausibility.py, MassPlausibility; data/statistical_mass.py.
% Fraction objectives use scale=max((high-low)/3.29,0.005).
% joint_upstream weak observation weight=0.01 in original-root-mass units.
% Integer step minimizes joint L1 displacement under the same constraints.
La riconciliazione determina insieme quantità lavorate, scarti e flussi in chilogrammi. Per ogni ODP candidato si sceglie una configurazione compatibile con gli intervalli delle quantità trasferite della \cref{sec:flow-selection-and-ambiguity}. Si mantengono i collegamenti scelti, aggiornandone le quantità. Al termine si salva questa sola configurazione, senza le alternative precedenti. L'univocità dei flussi salvati deriva quindi dalla scelta della procedura, non da nuove informazioni sul percorso fisico.

Per ogni rapporto stimato $\widehat{p}$, con intervallo empirico $[a,b]$, si usa la scala
\begin{equation}
\label{eq:mass-imputation-scale}
\tau_p=\max\left(\frac{b-a}{3.29},\,\num{0.005}\right).
\end{equation}
Il fattore $3.29\simeq2\cdot1.645$ è la larghezza dell'intervallo centrale al 90\% di una normale standard. La divisione fornisce una scala approssimata, poiché gli intervalli dei rapporti possono essere asimmetrici. Una scala più piccola dà maggior peso all'accordo con la stima.

Per ogni ODP candidato $k$, indichiamo con $\mathcal{R}_k$ la resa e gli scarti relativi parziali da riconciliare, e con $\mathcal{V}_k$ le quantità, gli scarti e i flussi da correggere. Ogni rapporto $p$ dipende dalle quantità da correggere, mentre la previsione $\widehat{p}$ e la scala $\tau_p$ restano fissate. Il simbolo $v^0$ indica il valore prima della riconciliazione e $q^0_{r(k)}$ la massa iniziale dell'ODP. L'obiettivo è
\begin{equation}
\label{eq:mass-reconciliation-objective}
\min_{\{v\}}\;
\sum_k\left[
\sum_{p\in\mathcal{R}_k}\left(\frac{p-\widehat{p}}{\tau_p}\right)^2
+\num{0.01}\cdot\sum_{v\in\mathcal{V}_k}
\left(\frac{v-v^0}{q^0_{r(k)}}\right)^2
\right].
\end{equation}
La procedura favorisce l'accordo con le stime, con una penalità debole sugli scostamenti dai valori precedenti. Il peso \num{0.01} è una scelta della procedura, non deriva da una stima dell'errore di misura. Le correzioni tendono quindi verso le previsioni, anziché fermarsi al primo valore entro le soglie.

La soluzione deve rispettare questi vincoli:
\begin{itemize}
\item Quantità finale uguale a quella registrata in \texttt{PRONTO}, \texttt{PRONTO\_QS} o \texttt{PRONTO\_QWR}, con bilanci esatti in ingresso e in uscita. Il bilancio terminale della \cref{sec:flow-mass-balances} ammette invece una tolleranza dello $0.1\%$.
\item Resa e scarti relativi valutabili entro i limiti MAD + IC95, mantenendo le modalità di scarto nullo, parziale e totale.
\item Almeno un chilogrammo su ogni collegamento scelto, con quantità lavorate e ripartizioni tra i rami entro i limiti descritti di seguito.
\end{itemize}

I riferimenti per questi ultimi limiti escludono i 90 ODV contenenti gli ODP candidati, come per i modelli della sezione precedente. Si applica il criterio MAD con incertezza bootstrap ai logaritmi delle quantità lavorate, dei loro rapporti con la massa iniziale e delle quote di flusso nelle diramazioni. Una quota è il flusso del collegamento diviso per il totale uscente dalla stessa occorrenza.

Per le occorrenze si cercano riferimenti con la stessa macrofase, forma, presenza di altri ODP nella posizione, modalità di scarto e possibilità di evitare il passaggio lungo un altro ramo. I gruppi vengono ampliati fino alla sola macrofase. Per le diramazioni si usano prima origine e destinazione, poi la sola origine. Le fasi composite includono nei riferimenti anche le macrofasi componenti. Servono almeno 20 osservazioni e limiti stimabili. Se mancano anche dopo l'ampliamento, le correzioni della posizione ODV non vengono applicate.

Il calcolo parte da un punto ammissibile che minimizza la somma degli scostamenti assoluti dai valori precedenti. Si ottimizza poi l'obiettivo con SLSQP, senza garanzia di ottimo globale, perché i rapporti dipendono dalle quantità da correggere. Infine, un problema lineare intero cerca quantità intere in chilogrammi quanto più vicine alla proposta, sempre in termini di scostamenti assoluti e rispettando tutti i vincoli.

Le correzioni degli ODP candidati della stessa posizione vengono applicate insieme soltanto se il calcolo e una verifica indipendente hanno esito positivo. Altrimenti si conservano i valori precedenti. Gli ODP non candidati, le quantità commerciali, le durate e i timestamp restano invariati.

\subsection{Effetti del trattamento}
\label{sec:mass-imputation-results}

Il trattamento viene applicato a tutti i 101 ODP candidati, distribuiti su 100 posizioni ODV e 90 ODV distinti. Delle 215 occorrenze materiali coinvolte, 193 cambiano valore e 22 rimangono invariate. La quantità in ingresso cambia in tutte le 193 occorrenze modificate e, in 152, cambia anche lo scarto. Le occorrenze modificate comprendono le 81 segnalate e altre 112 adeguate nella riconciliazione.

\begin{table}[htbp]
\centering
\small
\begin{tabularx}{\linewidth}{@{}Xrr@{}}
\toprule
\textbf{Grandezza} & \textbf{Prima} & \textbf{Dopo} \\
\midrule
Massa iniziale complessiva [kg] & \num{2357123} & \num{1970216} \\
Scarto complessivo [kg] & \num{486679} & \num{99629} \\
Materiale pronto registrato [kg] & \num{1870587} & \num{1870587} \\
Resa: percentile 5 [\%] & \num{47.15} & \num{93.27} \\
Resa: mediana [\%] & \num{80.74} & \num{94.72} \\
Resa: percentile 95 [\%] & \num{95.89} & \num{96.10} \\
\bottomrule
\end{tabularx}
\caption{Effetto della riconciliazione sui 101 ODP trattati. Le masse iniziali sono conteggiate una volta per ODP. Prima del trattamento i bilanci ammettono una tolleranza dello $0.1\%$, con scostamento massimo osservato di 34 kg nei 101 ODP. Dopo sono esatti in chilogrammi.}
\label{tab:mass-imputation-effects}
\end{table}


Sull'intero dataset, lo scarto complessivo passa da \num{8487.24} a \num{8100.19} t, con una riduzione del \num{4.56}\%. Il rapporto tra scarto totale e massa iniziale totale passa dal \num{5.31}\% al \num{5.08}\%, mentre la massa iniziale diminuisce dello \num{0.24}\%.

Massa iniziale e scarto complessivo diminuiscono in 95 ODP e aumentano in sei. Tutti i 76 casi segnalati per resa bassa appartengono al primo gruppo. La procedura considera affidabile il materiale pronto registrato e corregge le quantità a monte. Per esempio, in un ODP segnalato per scarto troppo basso, lo scarto passa da 1 a \num{2199} kg, con un aumento della massa iniziale da \num{41600} a \num{43799} kg.
% Case 000010884615: yield 100% -> 94.9793%; ready fixed at 41600 kg.
% Case 000010872293: the only treated ODP below 90% afterwards, at 89.6484%.
% Exact case values and predictions are in tables/data/statistical-mass-summary.json.

La \cref{fig:mass-imputation-effects} confronta i rapporti prima e dopo il trattamento con quelli dei \num{4759} ODP non trattati. Dopo la correzione, i percentili 5--95 della resa sono \num{93.27}--\num{96.10}\%, contro \num{90.49}--\num{97.03}\% nei non trattati. Per lo scarto relativo di \texttt{TRENOLUNA} sono \num{3.58}--\num{5.79}\%, contro \num{2.78}--\num{7.80}\%. L'imputazione concentra quindi i rapporti intorno alle previsioni e non riproduce tutta la variabilità osservata nei non trattati.

\begin{figure}[H]
\centering
\reportgraphic{statistical-mass-imputation.pdf}
\caption{Boxplot dei rapporti: le scatole vanno dal primo al terzo quartile, la linea interna indica la mediana e i baffi i percentili 5 e 95. I valori oltre i baffi non sono mostrati. Le prime due righe di ciascun pannello riguardano gli stessi 101 ODP prima e dopo il trattamento, la terza i \num{4759} ODP non trattati.}
\label{fig:mass-imputation-effects}
\end{figure}

% Real current-audit example: ODP 000010872769, ODV 0001116757, pos 000160.
% Statistical mass corrections: root 12518/438 -> 10449/486;
% fully scrapped BANCHI 2300/2300 -> 183/183; observed READY remains 9780.
% Trigger: global yield 9780/12518=78.1275%; local root scrap not flagged.
\begingroup
\tcbset{lines before break=12}
\begin{examplebox}[ex:joint-mass-imputation]{Una resa anomala e due occorrenze riconciliate}
Nell'ODP \texttt{000010872769}, la resa iniziale è $\num{9780}/\num{12518}\simeq\num{78.13}\%$, inferiore al limite MAD + IC95 di circa \num{84.54}\%. Lo scarto relativo alla quantità lavorata di \texttt{TRENOLUNA} non è segnalato. Un ramo ai \texttt{BANCHI} è interamente scartato e ha quindi $\rho_i=1$, escluso dal confronto statistico sugli scarti parziali. La riconciliazione conserva i \num{9780} kg registrati in \texttt{PRONTO} e i collegamenti della configurazione selezionata, modificando le due occorrenze materiali:

\begin{center}
\small
\begin{tabular}{@{}lrrrr@{}}
\toprule
& \multicolumn{2}{c}{\textbf{Prima [kg]}} & \multicolumn{2}{c}{\textbf{Dopo [kg]}}\\
\cmidrule(lr){2-3}\cmidrule(l){4-5}
\textbf{Macrofase} & \textbf{Ingresso} & \textbf{Scarto} & \textbf{Ingresso} & \textbf{Scarto}\\
\midrule
\texttt{TRENOLUNA} & \num{12518} & 438 & \num{10449} & 486\\
\texttt{BANCHI} & \num{2300} & \num{2300} & 183 & 183\\
\texttt{PRONTO} & \num{9780} & 0 & \num{9780} & 0\\
\bottomrule
\end{tabular}
\end{center}

Il modello prevede una resa del \num{94.34}\%, con intervallo empirico \num{90.35}--\num{96.48}\%. Dopo la riconciliazione, il bilancio è $\num{10449}-486-183=\num{9780}$ kg e la resa è circa il \num{93.60}\%: il risultato tiene conto sia delle stime sia dei vincoli e delle altre penalità. Il ramo ai \texttt{BANCHI} conserva lo scarto totale, ma cambia la quantità che lo percorre. Si modificano quindi anche dati non segnalati individualmente, perché la proposta deve soddisfare congiuntamente i vincoli statistici e di conservazione della massa.
\end{examplebox}
\endgroup

\section{Trattamento delle durate positive anomale}
\label{sec:positive-duration-treatment}

Il trattamento riguarda \texttt{durata\_totale\_minuti} e distingue le durate positive insolite da quelle incompatibili con il completamento del ciclo. Le prime vengono confrontate con altre durate della stessa macrofase. Per le seconde, già introdotte nella \cref{sec:temporal-duration-corrections}, servono anche previsioni e intervalli che delimitino le correzioni plausibili. In entrambi i casi, modificare una durata aggiorna la fine calcolata, senza cambiare l'inizio e la fine registrati.

\subsection{Individuazione delle durate anomale}
\label{sec:duration-outlier-detection}

% Sources: duration/statistical_audit.py; statistical_duration_imputation.json;
% expected_duration_imputation.json. Only the active MAD + IC95 criterion is shown.
% Counts and limits: statistical_duration_thresholds_260915.parquet and
% expected_duration_thresholds_260915.parquet. Evidence hashes in duration summary.
Si analizzano le durate positive di \texttt{BANCHI}, \texttt{CIRCO}, \texttt{TRATT\_LUNA} e \texttt{TRATT\_LFC}. Per ciascuna macrofase si applica il criterio MAD + IC95 della \cref{sec:mass-outlier-detection}, usando il logaritmo della durata al posto del logit. Le soglie tornano in minuti attraverso la funzione esponenziale. Restano gli stessi requisiti di numerosità e variabilità e i \num{2000} ricampionamenti per ODV. Le durate nulle non entrano nel confronto.

La \cref{tab:duration-outlier-thresholds} mostra i limiti e le segnalazioni. Sono individuate 12 durate troppo alte ai \texttt{BANCHI} e due troppo basse, una in \texttt{TRATT\_LUNA} e una in \texttt{TRATT\_LFC}. Non vi sono segnalazioni per \texttt{CIRCO}.

\begin{table}[htbp]
\centering
\small
\begin{tabularx}{\linewidth}{@{}Xrrrr@{}}
\toprule
\textbf{Macrofase} & $n$ & \textbf{Limite inf. [min]} & \textbf{Limite sup. [min]} & \textbf{Segnalati}\\
\midrule
\texttt{BANCHI} & \num{2368} & \num{0.80} & \num{2191.83} & 12 \\
\texttt{CIRCO} & \num{51} & \num{6.62} & \num{11943.22} & 0 \\
\texttt{TRATT\_LUNA} & \num{100} & \num{70.71} & \num{2292.34} & 1 \\
\texttt{TRATT\_LFC} & \num{187} & \num{194.63} & \num{19694.64} & 1 \\
\bottomrule
\end{tabularx}
\caption{Limiti MAD + IC95 per le durate positive. $n$ conta le occorrenze confrontate, prima delle sostituzioni. Per tutte e quattro le macrofasi sono valide le \num{2000} repliche bootstrap.}
\label{tab:duration-outlier-thresholds}
\end{table}


\subsection{Valori sostitutivi e controlli temporali}
\label{sec:duration-replacement-rules}

% Implementations: data/statistical_duration.py, clean_statistical_durations;
% data/expected_duration.py, clean_expected_durations; models/phase_duration.py.
% The mean rule recomputes admissible flows; TRATT_LFC keeps the saved graph.
% These are distinct from duration/correction_models.py used in Chapter 3.
Per \texttt{BANCHI}, \texttt{CIRCO} e \texttt{TRATT\_LUNA}, il valore proposto è la media delle durate positive non segnalate della stessa macrofase. Il riferimento è fissato prima delle sostituzioni e non viene aggiornato con i valori imputati. Le medie utilizzate sono \num{97.99} minuti per \texttt{BANCHI} e \num{452.42} per \texttt{TRATT\_LUNA}.

Per \texttt{TRATT\_LFC} si usa invece una regressione Ridge, con penalità $\lambda=10$, sul logaritmo della durata. Le variabili esplicative comprendono la quantità lavorata, la prima dimensione, la lunghezza richiesta, lo scarto relativo, il rapporto tra quantità lavorata e massa iniziale dell'ODP, prodotto, forma, finitura, stato fisico, marca dell'acciaio e classe di carbonio. Quantità lavorata, dimensione e lunghezza sono trasformate con il logaritmo. Il modello è addestrato su 154 occorrenze di 74 ODV, escludendo l'ODV del candidato.

Assumendo residui normali sulla scala logaritmica, si ottiene una durata con distribuzione lognormale. Indicando con $\mu_i$ la previsione del logaritmo e con $\sigma$ la scala dei residui, il valore atteso usato per la sostituzione è
\begin{equation}
\label{eq:duration-lfc-expected-value}
\widehat{d}_i=\exp\left(\mu_i+\frac{\sigma^2}{2}\right).
\end{equation}
La scala $\sigma$ è stimata dagli errori di previsione ottenuti con cinque suddivisioni per ODV, prevedendo ogni gruppo con un modello addestrato sugli altri quattro.

Le proposte basate sulla media devono rispettare i tempi di tutti i collegamenti già presenti. Si ricostruiscono poi i flussi e si richiedono un errore minimo di massa non peggiore di quello precedente e un'ambiguità entro la soglia ammessa. La nuova ricostruzione può semplificare il grafo. Le correzioni della stessa posizione ODV vengono applicate insieme solo se tutti i controlli riescono.

Per \texttt{TRATT\_LFC}, il grafo resta invece fissato. La durata prevista viene accettata se non introduce nuovi conflitti temporali sui collegamenti dell'ODP. Non viene accorciata per farla rientrare nel tempo disponibile. In caso di verifica negativa si mantiene il valore precedente. Entrambe le procedure conservano quantità, scarti e tempi registrati, aggiornando la fine calcolata e gli sfasamenti interessati.

\subsection{Modelli per le correzioni temporali}
\label{sec:duration-correction-models}

% Actual model used by current duration_corrections_260915.parquet:
% outputs/models/duration_corrections_v3/models.joblib, SHA-256
% 3d55fe973664ea66507c1bf32f232ea9a6c9cbcd3b1345fdbfa2ff5cfe87b08f.
% Training and selected metrics: analyze_ready_overrun_predictions.py and
% outputs/analysis/duration_correction_training_v3/.
% Training excludes 60 declared target OdV, but runtime target discovery is dynamic.
% Current applied correction ODP 000010902536 belongs to training OdV 0001126016.
% Do not claim that every runtime candidate has an out-of-OdV prediction.
Le correzioni della \cref{sec:temporal-duration-corrections} usano modelli distinti dalle sostituzioni appena descritte. Per ciascuna macrofase si prevede il logaritmo della durata in funzione di quantità lavorata, prima dimensione, lunghezza richiesta, prodotto, forma, finitura, stato fisico e marca dell'acciaio. La quantità è trasformata con \mbox{$\log(1+q_i)$}. I valori numerici mancanti sono sostituiti con le mediane dell'addestramento, mentre quelli categoriali formano una categoria dedicata.

I riferimenti hanno quantità e durate positive e appartengono a ODP senza occorrenze materiali oltre la chiusura del pronto. Esclusi i 60 ODV riservati ai casi da correggere durante la costruzione dei modelli, i riferimenti sono divisi per ODV in \num{1326} ODV di addestramento, 443 di calibrazione e 443 di verifica. Nell'addestramento si eliminano i logaritmi delle durate sotto il primo quartile o sopra il terzo di oltre \num{1.5} volte la distanza tra i due quartili.

Si confrontano Ridge, Random Forest e CatBoost \parencite{prokhorenkova2018catboost} mediante tre suddivisioni interne per ODV. Per ogni macrofase si sceglie il modello con il minore errore assoluto medio sulla scala logaritmica. La \cref{tab:duration-correction-validation} riporta i risultati dei modelli scelti su ODV di verifica separati. Se i riferimenti non bastano a costruire e calibrare un modello, non viene fornita una previsione.

\begin{table}[htbp]
\centering
\small
\setlength{\tabcolsep}{4pt}
\begin{tabularx}{\linewidth}{@{}Xlrrrr@{}}
\toprule
\textbf{Macrofase} & \textbf{Modello} & \textbf{Add. / cal.} & \textbf{Verifica} & \textbf{MAE [min]} & \textbf{Cop. [\%]}\\
\midrule
\texttt{ANTIMIX} & CatBoost & \num{884} / \num{321} & \num{244} & \num{16.58} & \num{93.03} \\
\texttt{BANCHI} & CatBoost & \num{2525} / \num{767} & \num{805} & \num{127.40} & \num{93.04} \\
\texttt{CONF} & Random Forest & \num{180} / \num{66} & \num{41} & \num{22.58} & \num{95.12} \\
\texttt{RADDR} & CatBoost & \num{550} / \num{206} & \num{151} & \num{91.49} & \num{92.05} \\
\texttt{TAG} & Ridge & \num{1089} / \num{354} & \num{411} & \num{67.11} & \num{93.67} \\
\bottomrule
\end{tabularx}
\caption{Verifica dei modelli usati nelle correzioni temporali del \cref{sec:data-and-treatment}. Le numerosità contano le occorrenze di addestramento, calibrazione e verifica. Il MAE è espresso in minuti e la copertura riguarda gli intervalli nominali al 90\%. La scelta del modello usa invece l'errore sulla scala logaritmica nelle suddivisioni interne all'addestramento.}
\label{tab:duration-correction-validation}
\end{table}


La calibrazione considera, per ciascun ODV, il maggiore errore assoluto tra logaritmo osservato e previsto. Per la macrofase $h$, con $G_h$ ODV di calibrazione, $c_h$ è il valore di rango $\lceil0.9\cdot(G_h+1)\rceil$ nell'elenco ordinato di questi errori. Indicando con $\widehat{\ell}_i$ il logaritmo previsto, la durata e il suo intervallo nominale al 90\% sono
\begin{equation}
\label{eq:duration-correction-interval}
\widehat{d}_i=\exp(\widehat{\ell}_i),
\qquad
I_i^{90}=\left[\exp(\widehat{\ell}_i-c_h),\,\exp(\widehat{\ell}_i+c_h)\right].
\end{equation}
Gli intervalli delimitano le durate candidate, mentre i controlli sul ciclo decidono se applicarle. La copertura misurata sui riferimenti non dimostra l'accuratezza delle correzioni. Inoltre, l'esclusione dei 60 ODV è fissata nell'addestramento e non viene ripetuta sui candidati individuati nell'applicazione: un ODV con una correzione applicata appartiene anche al gruppo di addestramento.

\subsection{Effetti del trattamento statistico}
\label{sec:duration-imputation-results}

% Counts here concern the positive-duration statistical stages only, not the
% completion-overrun corrections described in Chapter 3 or zero-duration imputation.
Le 14 durate segnalate dal confronto statistico appartengono a 14 ODP, in altrettante posizioni ODV e ODV distinti. Tutte le proposte superano i controlli. Le 12 durate dei \texttt{BANCHI}, comprese tra \num{2336.72} e \num{7862.10} minuti, diventano \num{97.99} minuti. Quella di \texttt{TRATT\_LUNA} passa da \num{3.32} a \num{452.42} minuti e quella di \texttt{TRATT\_LFC} da \num{12.96} a \num{1294.18} minuti.

La \cref{fig:duration-imputation-effects} confronta queste variazioni con le durate positive non segnalate. La sostituzione con la media assegna lo stesso valore a tutte le occorrenze trattate della macrofase, senza riprodurne la variabilità. Nei 13 ODP trattati con questa regola, la nuova ricostruzione riduce le diramazioni aggiuntive in sette casi e non richiede più il completamento implicito in quattro. Il trattamento di \texttt{TRATT\_LFC} conserva invece il grafo.

\begin{figure}[htbp]
\centering
\reportgraphic{statistical-duration-effects.pdf}
\caption{Durate positive corrette e riferimenti non segnalati della stessa macrofase. Le scatole mostrano primo quartile, mediana e terzo quartile, con baffi ai percentili 5 e 95. I cerchi indicano i valori precedenti e i rombi quelli sostitutivi, riportati anche in minuti. Le 12 sostituzioni dei \texttt{BANCHI} coincidono nello stesso rombo. L'asse orizzontale è logaritmico.}
\label{fig:duration-imputation-effects}
\end{figure}

% Current expected_duration_imputations audit, ODP 000010894433.
% Start 2026-05-01 05:47; next RADDR start 2026-05-03 09:15, cap 3088 min.
% Prediction 1294.1831536248073; registered end remains 2026-05-01 06:00.
\begingroup
\tcbset{lines before break=12}
\begin{examplebox}[ex:positive-duration-lfc]{Una durata troppo breve per il trattamento esterno}
Nell'ODP \texttt{000010894433}, \texttt{TRATT\_LFC} registra \num{12.96} minuti per \num{23375} kg. La durata è inferiore al limite MAD + IC95 di \num{194.63} minuti. Il modello propone \num{1294.18} minuti, circa 21 ore e 34 minuti, che rientrano nei \num{3088} minuti disponibili tra l'inizio del trattamento e l'inizio della successiva \texttt{RADDR}. La proposta supera i controlli sui collegamenti e viene applicata. Cambia la fine calcolata, mentre masse, grafo e tempi registrati restano invariati.
\end{examplebox}
\endgroup

\section{Imputazione delle durate nulle}
\label{sec:zero-duration-imputation}

Per \texttt{ANTIMIX}, \texttt{CONF} e \texttt{MOL} si assume che la lavorazione richieda un tempo positivo. Le durate registrate a zero in queste macrofasi vengono quindi sostituite. La regola riguarda queste tre macrofasi e non si estende agli altri zeri presenti nei consuntivi.

\subsection{Scelta dei valori sostitutivi}
\label{sec:zero-duration-models}

% Sources: data/zero_duration.py; data/rules/zero_duration_imputation.json;
% fit_additional_duration_distributions.py and its saved training metadata.
% Frozen model hashes and regenerated predictions: tables/data/zero-duration-summary.json.
Per \texttt{ANTIMIX} e \texttt{CONF} si estrae una durata da una distribuzione che dipende dalle caratteristiche dell'occorrenza. Si usano le stesse variabili esplicative del modello di \texttt{TRATT\_LFC} nella \cref{sec:duration-replacement-rules}, comprese quantità lavorata, dimensioni e caratteristiche del prodotto. Il logaritmo della durata è previsto con CatBoost per \texttt{ANTIMIX} e con Ridge, con $\lambda=10$, per \texttt{CONF}. Alla previsione si aggiunge un residuo casuale, normale nel primo caso e logistico nel secondo. Tornando ai minuti con la funzione esponenziale si ottengono, rispettivamente, una distribuzione lognormale e una log-logistica.

L'addestramento usa 663 durate positive di \texttt{ANTIMIX} e 79 di \texttt{CONF}, selezionate con il criterio MAD + IC95 calcolato sul solo gruppo di addestramento. Gli ODV sono separati tra addestramento e validazione. Le distribuzioni dei residui sono stimate dagli errori ottenuti con cinque suddivisioni per ODV, prevedendo ogni gruppo con un modello addestrato sugli altri quattro. I confronti tra modelli e distribuzioni usano la validazione, che quindi non costituisce una verifica indipendente delle scelte. I modelli sono fissati prima dell'imputazione e non vengono riaddestrati sui valori che producono.

Per \texttt{MOL} sono disponibili soltanto tre durate positive, pari a 60, 150 e 90 minuti. Si usa la loro media, 100 minuti, per tutti gli zeri. Questo riferimento limitato non permette di rappresentare la variabilità della lavorazione. Per le altre due macrofasi, invece, l'estrazione conserva una componente di variabilità anche a parità di caratteristiche. Le estrazioni sono riproducibili grazie a un seme casuale e a un ordinamento delle occorrenze fissati.

\subsection{Compatibilità con i tempi registrati}
\label{sec:zero-duration-time-windows}

Il grafo ricostruito determina il tempo disponibile per ciascuna occorrenza. Si considera il limite più vicino tra quelli imposti dai collegamenti in uscita e dal completamento dell'ODP, misurandolo dall'inizio registrato dell'occorrenza. Per \texttt{ANTIMIX} e \texttt{CONF}, quando questo tempo è positivo, si estrae dalla sola parte della distribuzione che vi rientra.

Indicando con $H_i$ la funzione di ripartizione della durata prevista e con $w_i>0$ il tempo disponibile, il valore estratto è
\begin{equation}
\label{eq:zero-duration-conditional-draw}
d_i^{\mathrm{imp}}=H_i^{-1}\!\left(u_i\cdot H_i(w_i)\right),
\qquad u_i\sim\mathcal{U}(0,1).
\end{equation}
La quantità $H_i(w_i)$ è la probabilità che una durata prevista rientri nel tempo disponibile. Se è inferiore al 5\%, il risultato viene segnalato come fortemente condizionato dai tempi: la sostituzione è applicata, ma proviene da una parte poco probabile della distribuzione. Per \texttt{MOL}, il valore di 100 minuti viene confrontato direttamente con il tempo disponibile.

Se non esiste un tempo positivo utilizzabile, oppure i 100 minuti di \texttt{MOL} non vi rientrano, la durata viene comunque imputata. Per \texttt{ANTIMIX} e \texttt{CONF} si usa un'estrazione senza il limite temporale, per \texttt{MOL} si mantengono 100 minuti. Questi casi sono marcati con \texttt{durata\_timestamp\_non\_vincolanti}: la procedura assume una durata positiva e segnala il conflitto con i tempi registrati, senza risolverlo. Il grafo, le masse e i timestamp registrati restano invariati. Vengono aggiornati la fine calcolata e gli sfasamenti.

\subsection{Risultati}
\label{sec:zero-duration-results}

% Counts verified against zero_duration_imputations and final consuntivi.
% All 1977 values regenerated from frozen models, saved windows and seed 42.
Sono sostituite \num{1977} durate nulle, appartenenti a \num{1901} ODP, \num{1786} posizioni ODV e \num{1169} ODV. La \cref{tab:zero-duration-results} distingue le macrofasi e i casi influenzati dai vincoli temporali. Le eccezioni sono 16: dieci \texttt{ANTIMIX} e tre \texttt{CONF} hanno tempo disponibile nullo, mentre tre \texttt{MOL} hanno soltanto 10, 15 e 37 minuti. Altre 87 estrazioni sono fortemente condizionate dai tempi disponibili.

\begin{table}[htbp]
\centering
\small
\begin{tabularx}{\linewidth}{@{}Xrrrr@{}}
\toprule
\textbf{Macrofase} & \textbf{Imputate} & \textbf{Mediana [min]} & \textbf{Condizionate} & \textbf{Eccezioni}\\
\midrule
\texttt{ANTIMIX} & \num{1940} & \num{14.78} & 86 & 10 \\
\texttt{CONF} & 23 & \num{40.86} & 1 & 3 \\
\texttt{MOL} & 14 & \num{100.00} & 0 & 3 \\
\bottomrule
\end{tabularx}
\caption{Sostituzione delle durate nulle. La colonna ``Condizionate'' conta i casi con forte condizionamento temporale, cioè con meno del 5\% della distribuzione prevista entro il tempo disponibile. Le ``Eccezioni'' sono i casi in cui tale tempo non viene imposto come vincolo. I due conteggi sono disgiunti.}
\label{tab:zero-duration-results}
\end{table}


La \cref{fig:zero-duration-distributions} confronta le durate imputate con quelle già positive prima di questo passaggio. Le mediane imputate di \texttt{ANTIMIX} e \texttt{CONF} sono inferiori a quelle dei rispettivi riferimenti, pari a 18 e 120 minuti. Il confronto descrive gruppi diversi per caratteristiche e vincoli temporali, e non misura l'accuratezza delle imputazioni: per le occorrenze registrate a zero la durata effettiva non è nota.

\begin{figure}[htbp]
\centering
\reportgraphic{zero-duration-distributions.pdf}
\caption{Durate già positive e durate imputate, con scatole dal primo al terzo quartile, mediana e baffi ai percentili 5 e 95. I riferimenti contengono 819 occorrenze di \texttt{ANTIMIX}, 94 di \texttt{CONF} e tre di \texttt{MOL}. Per \texttt{MOL}, le 14 imputazioni coincidono a 100 minuti. L'asse orizzontale è logaritmico.}
\label{fig:zero-duration-distributions}
\end{figure}

\begingroup
\tcbset{lines before break=12}
\begin{examplebox}[ex:zero-duration-window]{Una durata nulla con poco tempo disponibile}
Nell'ODP \texttt{000010874347}, un'occorrenza di \texttt{ANTIMIX} ha durata nulla e inizio registrato alle 15:49. Il materiale risulta pronto alle 16:02, lasciando 13 minuti disponibili. La distribuzione prevista assegna a questa finestra una probabilità di circa \num{0.55}\%. L'estrazione condizionata produce \num{12.21} minuti e porta la fine calcolata alle 16:01:12. La durata rientra nel tempo disponibile, ma viene segnalata come fortemente condizionata dal vincolo temporale.
\end{examplebox}
\endgroup
